\subsubsection{Visualisierung von Baumstrukturen}

Nach der Auswertung und Verarbeitung der LiDAR-Punktwolke soll eine adäquate Repräsentation der erfassten Baumstrukturen erfolgen. Hierfür existieren verschiedene Methoden, die sich im Wesentlichen in die Visualisierung direkter Punktwolken, die voxelbasierte Diskretisierung und die Modellierung mittels Polygonnetzen (Meshes) unterteilen lassen. Da diese Verfahren unterschiedliche geometrische Eigenschaften abbilden, können sie in der Praxis auch komplementär eingesetzt.

Die direkte Darstellung als LiDAR-Punktwolke nutzt die nativen 3D-Koordinaten der Messpunkte, wodurch die geometrischen Eigenschaften des Zielobjekts ohne Informationsverlust erhalten bleiben. da keine topologische Transformationen erforderliche sind, lassen sich die Daten unmittelbar übertragen und visualisieren. \citep{wangAutomatedPhenotypicTrait2024} Die Methode stößt jedoch bei stark verrauschten oder spärlich besetzten Datensätzen an informationstechnische Grenzen. Zudem ist die Verwaltung großflächiger Punktwolken rechenintensiv, was die Effizienz bei zunehmender Umgebungsgröße mindert. \citep{chengRealTimeLiDARInertial2024} Bei komplexen Strukturen können zudem Abschattungseffekte zu Verdeckungen und unvollständigen Datenreihen führen. Trotz dieser Limitierungen eignet sich die direkte Punktwolkendarstellung für die Erfassung feiner biologischer Merkmale und geometrischer Details, wie sie beispielsweise in der landwirtschaftlichen Phänotypisierung gefordert sind. Da die Technologie die vertikale Struktur präzise abbildet, lassen sich morphologische Parameter wie der Kronendurchmesser oder die Pflanzenhöhe direkt aus den diskreten Koordinaten ableiten. \citep{wangAutomatedPhenotypicTrait2024}Ein Ansatz mit höherem Vorverarbeitungsaufwand ist die Unterteilung des 3D-Raums in ein Gitter aus Rasterzellen (Voxel). Die Zuordnung der einzelnen LiDAR-Punkte zu diesen Voxeln schafft im Vergleich zu ungeordneten Punktmengen eine strukturierte Datenbasis. Zur Effizienzsteigerung speichern modernere Verfahren ausschließlich die explizit besetzten Voxel. \citep{chengRealTimeLiDARInertial2024} Bei klassischen voxelbasierten Methoden wächst der Speicher- und Rechenaufwand jedoch signifikant mit der Umgebungsgröße, da nachfolge Algorithmen das gesamte Gitter durchlaufen müssen. Zudem erweisen sich geometrische Berechnungen mittels Ray-Casting bei verrauschten oder lückenhaften LiDAR-Rohdaten als fehleranfällig. Die gewählte Voxelgröße beeinflusst maßgeblich die Darstellungsqualität, da eine zu grobe Auflösung zu visuellen Unschärfen und Verfälschungen der tatsächlicher Objektoberfläche führt. \citep{chengRealTimeLiDARInertial2024}

Die Modellierung mittels Meshes überführt Punkte zusammenhängender Oberflächen in geschlossene Dreiecksnetze. Zur Generierung dieser Netze existieren verschieden Ansätze. Zum einen lassen sich Oberflächen direkt aus voxelisierten Daten extrahieren. \citep{chengRealTimeLiDARInertial2024} Zum anderen ermöglichen geometrische Algorithmen wie Alpha-Shapes die Extraktion von Randpunkten, um über eine konstruierte Hüllkugel zusammenhängende Flächen zu berechnen. \citep{wangAutomatedPhenotypicTrait2024} Für hochkomplexe biologische Objekte wie Bäume kann Procedural Modelling eingesetzt werden. Hierbei werden zunächst Skelettstrukturen des Objekts rekonstruiert, auf denen anschließend das Mesh aufgebaut wird. \citep{keerthinathanModellingLiDARBasedVegetation2025} Eine zentrale Herausforderung liegt in ungeordneten und lückenhaften Punktwolken wie sie gerade bei Vegetation auftreten. Unabgetastete Bereiche führen zu Löchern im generierten Mesh-Modell und erzeugen Fehler bei der Oberflächenapproximation. \citep{chengRealTimeLiDARInertial2024} Zudem erfordert eine realistische Darstellung eine hohe Dreiecksdichte. Für präzise LiDAR-Sensorsimulationen wird beispielsweise eine Dichte von 11.000 Dreiecken pro Kubikmeter empfohlen, was den Rechen- und Speicheraufwand drastische erhöht. \citep{goodinGeometricFidelityRequirements2024}

Um die jeweiligen algorithmischen Schwächen auszugleichen, werden Voxel- und Mesh-Strukturen in vielen Anwendungen kombiniert. Während Voxel die volumetrische Innenstruktur effizient verwaltet, liefern Meshes die geometrische Außenhülle. In der Echtzeit-Navigation wird beispielsweise mithilfe von Voxeln eine performante Kollisionsabfrage im Raum realisiert, während hochauflösende Mesh-Karten die Geometrie der unbekannten Umgebung planen. \citep{chengRealTimeLiDARInertial2024} Ähnliche Synergien zeigen sich bei physikalischen Simulationen wie der Ausbreitung von Waldbränden. Voxel werden hierbei genutzt, um volumetrische Daten über die Verteilung von Brennstoffen sowie thermodynamische Eigenschaften innerhalb von Strukturen abzubilden. Das Mesh dient dabei als Begrenzungshülle, welche die Interaktion des Feuers mit der komplexen Topologie der Umgebung festlegt. \citep{goodinGeometricFidelityRequirements2024}

Für die Umsetzung in dieser Arbeit wurde die direkte punktwolkenbasierte Darstellung gewählt, da sie algorithmische Effizienz mit hoher visueller Detailtreue verbindet. Da LiDAR-Systeme Geometriedaten nativ erfassen, entfallen komplexe und rechenintensive Transformationen, wie sie bei Voxel- und Mesh-Generierung notwendig sind. Um den Rechenaufwand bei der Verarbeitung eines vollständigen Stadtgebiets zu begrenzen, beschränkt sich dieser Ansatz im Rahmen einer prototypischen Umsetzung auf die Visualisierung ausgewählter Abschnitte. Dies minimiert den Vorverarbeitungs- und Implementierungsaufwand erheblich. Zugleich garantiert dieser Ansatz für eine hohe visuelle Authentizität, da morphologische Details ohne geometrische Abstraktion erhalten bleiben. Unschärfe durch grobe Voxelauflösung oder fehlerhafte Oberflächenapproximationen werden auf diese Weise vermieden.